\documentclass[11pt]{amsart}
\usepackage{amsmath,amssymb,amsthm}
\usepackage[margin=1.15in]{geometry}
\usepackage[hidelinks]{hyperref}

\newtheorem{theorem}{Theorem}[section]
\newtheorem{lemma}[theorem]{Lemma}
\newtheorem{corollary}[theorem]{Corollary}
\theoremstyle{remark}
\newtheorem{remark}[theorem]{Remark}

\newcommand{\arxiv}[1]{\href{https://arxiv.org/abs/#1}{arXiv:\nolinkurl{#1}}}
\newcommand{\N}{\mathcal{N}}
\newcommand{\cl}[2]{N_{#1}[#2]}
\setlength{\emergencystretch}{2.5em}
\makeatletter\g@addto@macro\UrlBreaks{\do-}\makeatother

\begin{document}

\title[One-step nullhomotopy for \v{C}ech complexes]{A one-step nullhomotopy for \v{C}ech complexes of hypercube
graphs}
\author{Yitzchak Shmalo}
\date{26 November 2026}
\address{Einstein Institute of Mathematics, The Hebrew University of Jerusalem, Jerusalem, Israel}
\email{yspennstate@gmail.com}
\subjclass[2020]{55N31, 55U10, 05C12}
\keywords{\v{C}ech complex, hypercube graph, nerve, persistent homology, contiguity}

\begin{abstract}
For a graph $G$ and an integer $r\ge0$ let $\N(G,r)$ be the nerve of the closed balls of radius $r/2$ about the vertices
of $G$ in its geometric realization. Adams, Shukla and Singh proved that the inclusion of $\N(Q_n,r)$ into
$\N(Q_n,r+2)$ is null-homotopic for the hypercube graph $Q_n$, and asked whether the inclusion into $\N(Q_n,r+1)$ is
null-homotopic as well. We show that it is, for every $n\ge1$ and $r\ge0$. Consequently every bar in the reduced
\v{C}ech persistence barcode of $Q_n$ over integer scales lives at a single scale. The proof rests on a lemma for graph maps into a
bipartite graph that agree up to one edge at every vertex, and it applies equally to Cartesian products of trees.
\end{abstract}

\maketitle

\section{Introduction}

Let $G$ be a finite connected simple graph, and give its geometric realization the shortest path metric in which every
edge is isometric to $[0,1]$. For $r\ge0$ the \v{C}ech complex $\N(G,r)$ has vertex set $V(G)$, and a set of vertices is
a simplex when the closed balls of radius $r/2$ about them have a common point. Since $\N(G,r)=\N(G,\lfloor r\rfloor)$
for all real $r\ge0$~\cite[Section~2]{ASS}, the filtration is determined by its integer scales.

Adams, Shukla and Singh~\cite{ASS} studied these complexes for the hypercube graph $Q_n$, whose vertices are the words
in $\{0,1\}^n$ with edges between words at Hamming distance one. Their Theorem~6.1 states that the inclusion
$\N(Q_n,r)\hookrightarrow\N(Q_n,r+2)$ is homotopic to a constant map for all $n\ge1$ and $r\ge0$; the proof, a technique
they credit to \v{Z}.~Virk, shows that the maps induced by the coordinate projections
$p_i(x_1\dots x_n)=(0\dots0\,x_{i+1}\dots x_n)$ are contiguous into scale $r+2$. They show by the example $n=2$, $r=1$
that $\iota$ and the map induced by $p_1$ need not be contiguous as maps into $\N(Q_n,r+1)$, and
ask~\cite[Question~7.5]{ASS} whether the inclusions
$\N(Q_n,r)\hookrightarrow\N(Q_n,r+1)$ are also always null-homotopic. For Vietoris--Rips complexes of hypercube graphs
the corresponding one-step statement was proved by Adams and Virk~\cite{AV}, who expected that part of their work would
carry over to \v{C}ech complexes of hypercube graphs.

\begin{theorem}\label{thm:cube}
For every $n\ge1$ and every real $r\ge0$ the inclusion $\N(Q_n,r)\hookrightarrow\N(Q_n,r+1)$ is null-homotopic.
\end{theorem}

\begin{corollary}\label{cor:pers}
For every $n\ge1$, integer $r\ge0$, degree $q\ge0$ and abelian group $M$, the map
$\widetilde H_q(\N(Q_n,r);M)\to\widetilde H_q(\N(Q_n,r+1);M)$ induced by inclusion is zero. Over a field, every interval in
the barcode of the reduced persistent homology of $(\N(Q_n,r))_{r\in\mathbb Z_{\ge0}}$ consists of a single scale.
\end{corollary}

The complexes themselves need not be contractible: $\N(Q_2,1)$ is the $4$-cycle $Q_2$, whose first homology is
nonzero, so the barcodes in Corollary~\ref{cor:pers} are not empty in general.

The coordinate projections $p_i$ and $p_{i+1}$ already induce contiguous maps into $\N(Q_n,r+1)$ when $r$ is even; the
difficulty is at odd scales. There we pass from $p_i$ to $p_{i+1}$ through an auxiliary vertex map whose image lies in
one colour class of the bipartite graph $Q_n$. The argument uses only that the target is bipartite and that consecutive
maps differ by at most one edge at each vertex (Theorem~\ref{thm:main}). The same conclusion therefore holds for
Cartesian products of trees (Theorem~\ref{thm:trees}).

\section{Preliminaries}

For a vertex $v$ and an integer $k\ge0$ write $\cl kv=\{x\in V(G):d(x,v)\le k\}$, where $d$ is the graph distance.

\begin{lemma}[{\cite[Lemma~2.2]{ASS}}]\label{lem:ASS}
Let $G$ be a finite connected simple graph and $k\ge0$ an integer. The simplices of $\N(G,2k)$ are the subsets of the
sets $\cl kv$, $v\in V(G)$. The simplices of $\N(G,2k+1)$ are the vertices of $G$ and the subsets of the sets
$\cl kv\cup\cl kw$ over the edges $vw$ of $G$.
\end{lemma}

A \emph{graph map} $f\colon G\to H$ is a map $V(G)\to V(H)$ such that for every edge $uw$ of $G$ either $f(u)=f(w)$ or
$f(u)f(w)$ is an edge of $H$. A graph map does not increase distances, so it maps $\cl kv$ into $\cl k{f(v)}$. By
Lemma~\ref{lem:ASS} it maps $\cl kv\cup\cl kw$ for an edge $vw$ into a generating set of the same form if $f(v)\ne f(w)$,
and into $\cl k{f(v)}$ otherwise; the latter is a vertex or lies in $\cl k{f(v)}\cup\cl kb$ for a neighbour $b$ of
$f(v)$, so it is a simplex of $\N(H,2k+1)$. Hence $f$ induces a simplicial map $f_*\colon\N(G,r)\to\N(H,r)$ for every integer $r$.
For each $r$ write $\iota\colon\N(H,r)\hookrightarrow\N(H,r+1)$ for the inclusion.

Two simplicial maps $\phi,\psi\colon K\to L$ are \emph{contiguous} if $\phi(\sigma)\cup\psi(\sigma)$ is a simplex of $L$
for every simplex $\sigma$ of $K$. Contiguous maps have homotopic geometric realizations~\cite[Section~3.5]{Sp}: the straight-line homotopy between
them maps each simplex $\sigma$ into the closed simplex spanned by $\phi(\sigma)\cup\psi(\sigma)$.

\section{Maps that differ by at most one edge}

\begin{theorem}\label{thm:main}
Let $G$ and $H$ be finite connected simple graphs with $H$ bipartite, and let $f,g\colon G\to H$ be graph maps with
$d(f(v),g(v))\le1$ for every $v\in V(G)$. Then for every integer $r\ge0$ the maps
$\iota\circ f_*,\ \iota\circ g_*\colon\N(G,r)\to\N(H,r+1)$ are homotopic.
\end{theorem}

\begin{proof}
If $H$ has a single vertex then $f=g$. Otherwise every vertex of $H$ has a neighbour, since $H$ is connected.

\emph{Even scales, $r=2k$.} Let $\sigma$ be a simplex of $\N(G,2k)$. By Lemma~\ref{lem:ASS},
$\sigma\subseteq\cl kc$ for some vertex $c$, so $f(\sigma)\subseteq\cl k{f(c)}$ and $g(\sigma)\subseteq\cl k{g(c)}$. If
$f(c)g(c)$ is an edge then $f(\sigma)\cup g(\sigma)\subseteq\cl k{f(c)}\cup\cl k{g(c)}$, a simplex of $\N(H,2k+1)$. If
$f(c)=g(c)$ the union lies in $\cl k{f(c)}\cup\cl kb$ for any neighbour $b$ of $f(c)$. So $\iota f_*$ and $\iota g_*$
are contiguous. Bipartiteness is not used at even scales.

\emph{Odd scales, $r=2k+1$.} Let $V(H)=A\sqcup B$ be a bipartition. For $x\in V(H)$ the parity of $d(x,a)$ is the same
for all $a\in A$: it is even if $x\in A$ and odd if $x\in B$. Define $h\colon V(G)\to A$ by
\[
h(v)=\begin{cases}
\text{the one of }f(v),g(v)\text{ that lies in }A,&\text{if }f(v)\ne g(v),\\
f(v),&\text{if }f(v)=g(v)\in A,\\
\text{a fixed neighbour of }f(v),&\text{if }f(v)=g(v)\in B.
\end{cases}
\]
Since $f(v)$ and $g(v)$ are equal or adjacent, $h(v)\in A$ is at distance at most one from both. We show that for every
simplex $\sigma$ of $\N(G,2k+1)$ the set $f(\sigma)\cup h(\sigma)$ lies in $\cl{k+1}p$ for some vertex $p$, and is
therefore a simplex of $\N(H,2k+2)$ by Lemma~\ref{lem:ASS}.

A single vertex $\sigma=\{v\}$ is clear, with $p=f(v)$. Otherwise $\sigma\subseteq\cl ku\cup\cl kw$ for an edge $uw$ of
$G$, and $f(\sigma)\subseteq\cl k{f(u)}\cup\cl k{f(w)}$. If $f(u)=f(w)=c$, then $f(\sigma)\subseteq\cl kc$ and
$h(\sigma)\subseteq\cl{k+1}c$; take $p=c$. If $f(u)f(w)$ is an edge, one endpoint lies in $A$ and the other in $B$, and
we choose $p\in\{f(u),f(w)\}$ such that $d(p,a)\equiv k+1\pmod 2$ for all $a\in A$. Every point of $f(\sigma)$ is within
$k$ of one endpoint, hence within $k+1$ of $p$. For $v\in\sigma$,
\[
d(p,h(v))\le d(p,f(v))+d(f(v),h(v))\le k+2,
\]
and $d(p,h(v))\equiv k+1\pmod2$ because $h(v)\in A$. Hence $d(p,h(v))\le k+1$.

In particular $h(\sigma)$ is a simplex, so $h$ defines a simplicial map $\N(G,2k+1)\to\N(H,2k+2)$, and it is
contiguous to $\iota f_*$. The argument used only that $f$ is a graph map and that $h(v)\in A$ lies within one of $f(v)$,
so it applies equally to $g$. Thus $\iota f_*\simeq h\simeq\iota g_*$.
\end{proof}

\begin{corollary}\label{cor:chain}
Let $H$ be a finite connected bipartite graph, and suppose there are graph maps $f_0,f_1,\dots,f_m\colon H\to H$ with
$f_0$ the identity, $f_m$ constant, and $d(f_{j-1}(v),f_j(v))\le1$ for all $j$ and $v$. Then for every real $r\ge0$ the
inclusion $\N(H,r)\hookrightarrow\N(H,r+1)$ is null-homotopic.
\end{corollary}

\begin{proof}
Since $\N(H,r)=\N(H,\lfloor r\rfloor)$ and $\lfloor r+1\rfloor=\lfloor r\rfloor+1$, we may take $r$ to be an integer.
The inclusion equals $\iota\circ(f_0)_*$, and $\iota\circ(f_m)_*$ is constant. Theorem~\ref{thm:main} gives
$\iota(f_{j-1})_*\simeq\iota(f_j)_*$ for each $j$.
\end{proof}

\begin{proof}[Proof of Theorem~\ref{thm:cube} and Corollary~\ref{cor:pers}]
The graph $Q_n$ is bipartite, with the colour of a word given by the parity of its number of ones. Put $f_0$ the
identity and $f_j=p_j$ for $1\le j\le n$. Each $p_j$ maps an edge to an edge or to a single vertex, $p_n$ is constant,
and $p_{j-1}(x)$ and $p_j(x)$ differ at most in coordinate $j$. Corollary~\ref{cor:chain} gives
Theorem~\ref{thm:cube}. A null-homotopic map induces zero on reduced homology with any coefficients, which is
Corollary~\ref{cor:pers}; over a field the interval decomposition of the persistence module then has no interval
containing two consecutive scales.
\end{proof}

\begin{remark}
At even scales Theorem~\ref{thm:main} shows that $\iota(p_{j-1})_*$ and $\iota(p_j)_*$ are contiguous. At odd scales the
intermediate map is needed. For instance, in $Q_2$ at scale $r=1$, the identity and $p_1$ are not contiguous into
$\N(Q_2,2)$, which is the boundary of a tetrahedron on the four vertices~\cite[Section~6]{ASS}. The
map $h$ of the proof sends all four vertices into the colour class $\{00,11\}$, and both contiguities hold.
\end{remark}

\section{Products of trees}

\begin{theorem}\label{thm:trees}
Let $G=T_1\,\square\,\cdots\,\square\,T_m$ be a Cartesian product of finite trees. Then for every real $r\ge0$ the
inclusion $\N(G,r)\hookrightarrow\N(G,r+1)$ is null-homotopic.
\end{theorem}

\begin{proof}
A Cartesian product of bipartite graphs is bipartite. For a single tree $T$ with $t\ge2$ vertices, let
$T=S_0\supset S_1\supset\dots\supset S_{t-1}$ be subtrees with $S_{i}=S_{i-1}\setminus\{\ell_i\}$ for a leaf $\ell_i$ of
$S_{i-1}$ with neighbour $b_i$ in $S_{i-1}$. Let $\rho_i\colon S_{i-1}\to S_i$ fix every vertex except $\ell_i$, which
it sends to $b_i$, and set $\phi_0=\mathrm{id}_T$ and $\phi_i=\rho_i\circ\phi_{i-1}$. Each $\phi_i$ is a graph map
$T\to T$, $\phi_{t-1}$ is constant, and $\phi_{i-1}(x)$ and $\phi_i(x)$ are equal or adjacent for every $x$.

Write $t_s=|V(T_s)|$. In $G$, run through these sequences one factor at a time: first apply $\phi_1,\dots,\phi_{t_1-1}$ to the first
coordinate with the others fixed, then the sequence for $T_s$ to coordinate $s$ for $s=2,\dots,m$ in turn. An edge of $G$ changes
one coordinate along an edge of that factor, and its image changes at most that coordinate along an edge or not at all,
so every map in the sequence is a graph map of $G$. Consecutive maps differ in one coordinate by at most one edge, and
the last map is constant. Corollary~\ref{cor:chain} applies.
\end{proof}

Theorem~\ref{thm:cube} is the case in which every factor is the edge $K_2$.

\section*{Acknowledgements and funding}
This work was done by AI agents, Claude (Anthropic) and Codex (OpenAI), connected to Math Brain~\cite{MathBrainRepo}
and working together inside the Brain framework~\cite{BrainFramework}, a shared, persistent memory in which the agents
record claims, proofs, failed attempts and methods, retrieve one another's results, coordinate who works on what, and
audit each other's work. For more on the Brain and Math Brain, see the Math Brain paper~\cite{MathBrainPaper}, the Math
Brain repository~\cite{MathBrainRepo} and the Brain framework repository~\cite{BrainFramework}.

The research presented in this paper was supported by the European Research Council (ERC) under the European Union's
Horizon Europe research and innovation programme (grant agreement No.~101041711), by the Simons Foundation as part of
the Collaboration on the Mathematical and Scientific Foundations of Deep Learning, by Heights Labs, by Convex Nexus
Capital, by the Israel Science Foundation (grant number 2258/19), and by the Israel Science Foundation (ISF Grant
4101/25).

\section*{Declaration of competing interest}
The author holds equity in Heights Labs and in Convex Nexus Capital and has a management role at Convex Nexus Capital.
These entities are included in the funding declaration above.

\section*{Data availability}
No data were generated or analysed in this work. A Python script that tests both contiguities of the proof of
Theorem~\ref{thm:cube} for $n\le5$ and all $r\le2n$, computing the \v{C}ech complexes from points of the geometric
realization, is at \url{https://github.com/yspennstate/cech-hypercube-one-step/tree/v1.1}; the command
\texttt{python check\_one\_step.py 5} runs it for $n\le5$.

\section*{Declaration of generative AI and AI-assisted technologies}
This paper is part of a research project on AI proving problems. Codex agents (OpenAI) produced the proof, and Claude
agents (Anthropic) checked the proof and the references and wrote the manuscript. The question is due to Adams, Shukla
and Singh~\cite{ASS}, and the coordinate projections and the use of contiguity are taken from their proof of
Theorem~6.1. The corresponding one-step statement for Vietoris--Rips complexes is due to Adams and Virk~\cite{AV}. The
author directed the work and is responsible for the content of the paper and for decisions concerning
its dissemination.

\begin{thebibliography}{9}
\bibitem{ASS} H.~Adams, S.~Shukla and A.~Singh, \emph{\v{C}ech complexes of hypercube graphs}, Homology Homotopy Appl.
\textbf{27} (2025), no.~1, 83--105, \href{https://doi.org/10.4310/HHA.2025.v27.n1.a6}{doi:10.4310/HHA.2025.v27.n1.a6};
\arxiv{2212.05871}. Numbered results are cited as in the arXiv version.
\bibitem{AV} H.~Adams and \v{Z}.~Virk, \emph{Lower bounds on the homology of Vietoris--Rips complexes of hypercube
graphs}, Bull. Malays. Math. Sci. Soc. \textbf{47} (2024), no.~3, Paper No.~72,
\href{https://doi.org/10.1007/s40840-024-01663-x}{doi:10.1007/s40840-024-01663-x}; \arxiv{2309.06222}.
\bibitem{Sp} E.~H.~Spanier, \emph{Algebraic Topology}, McGraw-Hill, New York, 1966.
\bibitem{MathBrainRepo} Y.~Shmalo, \emph{Math Brain: public snapshot, code and tests}, 2026,
\url{https://github.com/yspennstate/math-brain-public/tree/5427734deee9be643d6dfbdd5e7d4baeeefe4a4a}.
\bibitem{MathBrainPaper} Y.~Shmalo, \emph{Math Brain: shared, auditable memory for multi-agent mathematical
research}, 2026, \href{https://doi.org/10.5281/zenodo.24150837}{doi:10.5281/zenodo.24150837},
\url{https://github.com/yspennstate/math-brain-public/tree/5427734deee9be643d6dfbdd5e7d4baeeefe4a4a/papers/math-brain-2026-10-24}.
\bibitem{BrainFramework} Y.~Shmalo, \emph{Brain Framework}, public software repository, 2026,
\url{https://github.com/yspennstate/brain-framework-public/tree/46b935d26d5a2893f7790b60466387dd8da47ef5}.
\end{thebibliography}

\end{document}
